\documentclass[aps,floats,showpacs,amssymb,tightenlines]{revtex4}



\usepackage{amsmath}
\usepackage{amsfonts}
\usepackage{amscd}
\usepackage{epsfig}
\usepackage{amssymb}
\usepackage{tabularx}
\def\a{\alpha}
\def\b{\beta}
\def\g{\gamma}
\def\G{\Gamma}
\def\e{\epsilon}
\def\m{\mu}
\def\n{\nu}
\def\d{\delta}
\def\l{\lambda}
\def\L{\Lambda}
\def\t{\theta}
\def\w{\omega}
\def\c{\nabla}
\def\p{\partial}
\def\k{\kappa}
\def\r{\rho}
\def\P{{\cal P}}
\def\I{{\cal I}}
\def\re{r_{ext}}
\def\ri{r_{int}}
\def\A{{\cal A}}
\def\po{\psi_{\w_o}}
\def\pe{\psi_{\w_o}^{(j=1)}{}}



\begin{document}
\title{Exciting the unstable modes of a negative mass Schwarzschild spacetime}

\author{Gustavo Dotti and Reinaldo J. Gleiser}
\affiliation {Facultad de Matem\'atica, Astronom\'{\i}a y
F\'{\i}sica (FaMAF), Universidad Nacional de C\'ordoba. Ciudad
Universitaria, (5000) C\'ordoba,\ Argentina}


\begin{abstract}
Schwarzschild spacetime with negative mass $M$ is known to admit gravitational instabilities
in the form of exponentially growing solutions
of Zerilli's equation, satisfying appropriate boundary conditions.
However, it is not known how to evolve general initial data using Zerilli's equation.
The reason is that when $M<0$ the Zerilli function, as defined in terms of regular metric perturbation,
has in general a simple pole at $r=r_s>0$, and also the
 Zerilli potential has
a second order pole at $r_s$. As a consequence, Zerilli functions are not square integrable,
and it is not possible to regard the  radial piece
 of Zerilli's equation as  a self-adjoint  operator, whose mode
expansion allows
the evolution of arbitrary initial data, as is done in the positive mass case. We do not know
 how to evolve
initial data even after having specified  boundary conditions at the $r=0$ singularity. In this paper we show
how this problem can be overcome using the technique of intertwining potentials. This allows
 to fill in some gaps in the proof of instability of the negative mass Schwarzschild spacetime,
 by showing that generic initial
data supported away from the $r=0$ singularity do in fact excite the unstable modes.
 \end{abstract}


\pacs{04.50.+h,04.20.-q,04.70.-s, 04.30.-w}

\maketitle

\section{Introduction}
The linear stability under gravitational perturbations of the negative mass Schwarzschild
spacetime was first considered in \cite{ghi}, where a proof of stability for the
vector (or odd) modes is given. The scalar (even) modes, reconsidered in \cite{dg},  are far more subtle,
because the behaviour of the Zerilli potential $V_z$ \cite{zerilli,ki} at the $r=0$ singularity implies
a one parameter ambiguity in boundary conditions at this point (parameterized by $S^1$ \cite{ghi}).
 The potential $V_z$ also has
a second order pole at $r=r_s$ ($r_s > 0$) in the domain of interest, referred to
as a ``kinematic'' singularity in \cite{dg}. None of these problems
are present in the positive mass case, for which $r > 2M$ and the kinematic singularity occurs for negative $r$.

The ambiguity in boundary conditions at $r=0$ was addressed to in \cite{ghi,dg} where it was shown  that, requiring
that the first order corrections to the Riemann tensor algebraic invariants do not diverge faster, as $r \to 0^+$, than
their zeroth order piece,
selects a unique boundary condition. Note that this requirement is natural if we want our first order
formalism to provide approximate solution of Einstein's equations that can be interpreted as arbitrarily small perturbations of the unperturbed metric. Also, this same choice of boundary conditions at
$r=0$
selects perturbations with finite energy, using the energy notion obtained by going to
second order perturbation theory \cite{ghi}.   Regarding
the kinematic singularity, its origin lies in the fact that, as defined, the Zerilli function $\Psi_z$ has a built in simple pole
at this point for generic smooth gravitational perturbations \cite{dg}. The Zerilli function is still useful in
decoupling the full set of linearized Einstein's equations in the negative mass case, but the usual
approach of solving Zerilli's equation  by regarding its radial piece as a self adjoint operator in $L^2((0,\infty),dx)$,
$dx=dr / (1-2M/r)$
breaks down, since $\Psi_z$ does not belong to this space, and $V_z$ has a second order pole
at $r=r_s>0$. This problem cannot be fixed by introducing alternative radial variables or integrating factors, which
can be seen to merely move the singularity from the coefficients of the differential equation to the measure
that makes its radial piece self adjoint. Since we lack a self adjoint radial operator in the space of physically relevant
Zerilli functions, we do not actually know how to evolve initial data $\Psi_z(t=0,r), \p \Psi_z / \p t (t=0,r)$, even after having specified boundary conditions at $r=0$.

In \cite{dg}, unstable (exponentially growing in time) solutions of Zerilli's equation satisfying the chosen
boundary condition at $r=0$, exponentially decaying for large $r$,
and leading to a smooth perturbed metric, are  given in an explicit form. It is also argued in \cite{dg}
that these unstable modes can be excited by initial data  compactly supported
away from $r=0$. Yet, since we do not know how to evolve in time a given initial configuration, we do not
have a proof of the fact that these modes can actually be excited by such  initial data.

In this paper we show how the problem of evolution of initial data for gravitational scalar
perturbations on a negative mass Schwarzschild spacetime can be solved using the technique of intertwining potentials
 (see \cite{price} and references therein).
An intertwining operator is constructed that sends physically relevant Zerilli functions onto $L^2((0,\infty),dx)$ and that,
at the same time, converts the Zerilli potential into another potential, free of kinematic singularities and having a unique
self adjoint extension that corresponds precisely to our physically motivated choice of boundary condition at $r=0$.
In Section \ref{rw} we give a brief account of Zerilli's approach to (scalar type) gravitational
perturbations of Schwarzschild spacetime,
stressing the problems that arise
 in the negative mass case. In Section \ref{inter} we
 introduce the intertwining operator technique. An account of the results we need from
 \cite{price} is given in Lemma 1, then we prove some further
results that are relevant to our problem (Lemma 2 and the Theorem below).
In Section \ref{apli} the intertwining technique is applied to Zerilli's equation for a negative mass
Schwarzschild spacetime, and shown to map the space of physically relevant Zerilli functions onto $L^2((0,\infty),dx)$,
while converting  Zerilli's potential into one free of kinematic singularities. In the  construction of  the
intertwining operator, the generalization to higher harmonics of the zero mode found in \cite{ghi} is obtained and used.
Section \ref{sum} summarizes our results.




\section{Scalar perturbations of the negative mass Schwarzschild spacetime}\label{rw}
  In the Regge-Wheeler gauge \cite{rw}, the scalar
perturbations for the angular mode $(\ell,m)$  are
described by  four functions $H_0(r,t)$, $H_1(r,t)$, $H_2(r,t)$ and
$K(r,t)$, in terms of which the perturbed metric takes the form,
\begin{eqnarray}
\label{RW1}
ds^2 & = & -\left(1-\frac{2M}{r}\right)\left(1- \e H_0 Y_{\ell,m}\right)dt^2 +
2 \: \e \;  H_1 Y_{\ell,m} dt dr+\left(1-\frac{2M}{r}\right)^{-1}\left(1+ \e H_2 Y_{\ell,m}\right)dr^2
\nonumber \\
& & +r^2\left(1+ \e K\, Y_{\ell,m}\right) \left( d \theta^2 +  \sin^2(\theta)
d \phi^2 \right)
\end{eqnarray}
where $\epsilon$ is an auxiliary parameter, used to  keep track of the order of the perturbation, and
$Y_{\ell,m}=  Y_{\ell,m}(\theta,\phi)$ are standard spherical
harmonics on the sphere. The linearized Einstein equations for the
metric (\ref{RW1}) imply $H_0(r,t)=H_2(r,t)$, and a set of coupled
differential equations for $H_1$, $H_2$ and $K$, which can be  ``diagonalized" by introducing
the Zerilli function $\Psi_z(t,r)$, by the replacements,
\begin{eqnarray}
\label{RW2}
K & = & q(r) \Psi_z+\left(1-\frac{2M}{r}\right)\frac{\partial \Psi_z}{\partial r} \nonumber \\
H_1 & = & h(r)\frac{\partial \Psi_z}{\partial t}+r\frac{\partial^2 \Psi_z}{\partial t \partial r}
 \\
H_2 & = & \frac{\partial }{\partial r}\left[\left(1-\frac{2M}{r}\right)\left(h(r) \Psi_z
+r\frac{\partial \Psi_z}{ \partial r} \right)\right] -K  \nonumber
\end{eqnarray}
where,
\begin{eqnarray}
\label{RW2a}
q(r) & = &  \frac{\lambda(\lambda+1)r^2+3 \lambda M r + 6 M^2}{r^2(\lambda r + 3 M)}\nonumber \\
h(r) & = &  \frac{\lambda r^2 - 3 \lambda r M -3 M^2}{(r-2M)(\lambda r + 3M)}\nonumber
\end{eqnarray}
and
\begin{equation}
\lambda =  \frac{(\ell-1)(\ell+2)}{2} .
\end{equation}
The linearized Einstein's equations then reduce to Zerilli's wave equation
\begin{equation}
\label{zerilli}
\frac{\partial^2 \Psi_z}{\partial t^2}+{\cal H}_z \Psi_z = 0
\end{equation}
where,
\begin{equation} \label{scho}
{\cal H}_z = -\frac{\partial^2 }{\partial x^2} +V
\end{equation}
looks like a  quantum Hamiltonian operator  with potential
\begin{equation}
\label{poten1}
V = 2 \left(1 -\frac{2 M}{ r}\right) {\frac{ \lambda^2 r^2
\left[(\lambda+1) r + 3 M \right] + 9 M^2 (\lambda r +M)}{ r^3
(\lambda r + 3 M)^2} } ,
\end{equation}
and $x$ is the ``tortoise'' coordinate, related to $r$ by $\frac{dx}{dr}=\left(1-\frac{2M}{r}\right)^{-1}$.
We choose the integration constant such that $x=0$ for $r=0$, then
\begin{equation}\label{tortu}
x = r +2 M \ln \left| \frac{r-2M}{2M}\right|.
\end{equation}
The usual approach to (\ref{zerilli}) is separation of variables:  $e^{\pm i E^{1/2} t} \psi_E(x)$
is a solution of (\ref{zerilli}) if ${\cal H}_z \psi_E = E \psi_E$.
In the positive mass case,   the exterior static region
$r>2M$ of a Schwarzschild black hole gets mapped onto $-\infty < x< \infty$,
the black hole horizon sitting at $x=-\infty$,
and  the relevant solutions
of (\ref{zerilli})
belong to $L^2({\mathbb R},dx)$, where ${\cal H}_z$ is self adjoint, as follows
from (\ref{scho}) and (\ref{poten1}).
In this case one solves (\ref{zerilli}) by expanding $\Psi_z$ and $\p \Psi_z / \p t$ at $t=0$
 in a complete set $\psi_E$ of eigenfunctions of ${\cal H}_z$,   then (\ref{zerilli}) reduces
 to  the following ordinary differential equations for $a_E(t) \equiv \int \psi_E(x) ^* \Psi_z(t,x) dx$:
\begin{eqnarray}
 \ddot a_E &=& -E a_E \\
 \dot a_E(0) &=& \int \psi_E ^* \left.  \left(\frac{\p \Psi_z}{ \p t}\right)
 \right \vert _{t=0}  \; dx  \\
 a_E(0) &=& \int \psi_E ^* \Psi_z |_{t=0} \; dx
\end{eqnarray}
  Moreover, since ${\cal H}_z$ is positive definite in this case, the  $E$ eigenvalues  are positive, $e^{\pm i E^{1/2} t} \psi_E(r)$
  oscillatory, and $\Psi_z$
   bounded at all times,
    for any initial conditions. This proves that the exterior, static region of a Schwarzschild black hole is stable.\\
In the negative mass case, the domain of (\ref{zerilli}) is $x>0$, so we need to specify boundary conditions at
 $x=0$ besides the initial values of $\Psi_z$ and $\p\Psi_z / \p t$, in order that (\ref{zerilli}) has  a unique solution.
 Also, the potential has a singularity at the boundary,
 \begin{equation} \label{vcero}
 V \simeq -1/(4x^2)+ ... \;\; \text{ for } \; x \to 0^+.
 \end{equation}
 Consider a quantum mechanical problem ${\cal H} \psi_E = E \psi_E$ on the half axis $x>0$, with a smooth potential
 behaving as (\ref{vcero})  and vanishing at infinity.
 For {\em any} $E$, the general local solution for $x \gtrsim 0$
  of the differential equation ${\cal H} \psi_E = E \psi_E$ is
 \begin{equation} \label{cero}
 \psi_E \simeq a_0 \left[ \left(\frac{x}{|M|}\right)^{1/2}  + ... \right] + b_0 \left[  \left(\frac{x}{|M|}\right)^{1/2}  \ln \left(\frac{x}{|M|} \right) + ... \right]
 , \;\;\; x \gtrsim 0.
 \end{equation}
 with  $E$ showing up at higher order.
Note that (\ref{cero}) is
 square integrable near $x=0$ for any $a_0$ and $b_0$, a peculiarity of  the potential
 $V  \simeq \a/x^2$ when $\a < 3/4$ (for $\a \geq 3/4$, only one of the local solutions is square
 integrable near $x=0$ \cite{meetz}.) If $E$ is negative,  there is always an asymptotic solution $\psi_E^o$ which is square
 integrable at infinity. Since $\psi_E^o$  behaves as (\ref{cero}) near zero, then $\psi_E^o \in L^2((0,\infty),dx)$.
 Such a potential allows a square integrable eigenfunction for any negative $E$!
 This introduces an ambiguity, as a choice  of an appropriate
 boundary condition needs to be imposed  to define a domain where ${\cal H}_z$ is self adjoint.
 The allowed choices for this purpose can be  parameterized by $S^1$, with  $(a_0,b_0) \propto (\cos \theta, \sin \theta)$
 \cite{ghi,meetz,reed}. In the space $S_{\theta}$ of functions behaving like (\ref{cero}) with
 $(a_0,b_0) \propto (\cos \theta, \sin \theta)$,
 the continuum of negative eigenvalues get  discretized, ${\cal H}$ is self adjoint, and arbitrary functions can be expanded in a basis of ${\cal H}$ eigenfunctions \cite{ghi,meetz,reed}.

However, as discussed in Section 7 of \cite{dg}, the quantum mechanical problem associated to ${\cal H}_z$ is not directly relevant to the gravitational perturbation problem when $M<0$.
To understand this crucial difference between the positive and negative mass cases,
consider the  relations (\ref{RW2}), which can be inverted and give
\begin{equation}
\label{zeri}
\Psi_z(r,t)  = \frac {r (r-2M)}{ ( \lambda+1 )  ( \lambda
\,r+3\,M ) } \left(H_2
 - r  \frac {\partial K}{\partial r}
\right)
    +\frac{r}{\lambda+1} K.
\end{equation}
Zerilli's function
succeeds in reducing the full set of linearized Einstein's equations to a single wave equation, however, for $M<0$, this function
is singular in the relevant $r>0$ range. The ``kinematic'' singularity at $r_s=-3M/\lambda$ in (\ref{zeri}) indicates
that physically acceptable Zerilli functions, i.e., those corresponding to smooth metric perturbations decaying for large $r$,
 have a simple pole at $r_s$. This built in singularity in $\Psi_z$ is certainly related  to the singularity in the
coefficients of the wave equation it satisfies, the second order pole at $r_s$ in the potential (\ref{poten1}).
Thus, even if  ${\cal H}_z$ were extended to a  self adjoint operator in some subspace of  $L^2({\mathbb R},dx)$,
this space would not be the natural setting for
physically acceptable gravitational perturbations, which, because of the kinematic singularity, generically, are not square integrable. In fact, the relevant Zerilli functions
are those admitting a Laurent expansion
\begin{equation} \label{crc}
\Psi_z = \sum_{j \geq -1}  c_j (r-r_s)^j, \;\;\; c_0 = \frac{\lambda^2  c_{-1}}{3M (3 + 2 \lambda)}
\end{equation}
for $r$ near $r_s$. The second of the above equations (which fixes a sign error in the analogous equation (16) in \cite{dg})
follows from requiring that the metric functions in (\ref{RW2}) be $C^2$ at $r_s$.

Regarding the boundary condition at $x=0$, we note that (\ref{cero}) gives the behaviour of any
solution of (\ref{zerilli})  in the $x =0$ limit (the $t$ dependence showing up at higher orders). As proved in \cite{ghi,dg}, if we did not make the choice $b_0=0$, the first order correction
to the Kretschmann invariant would diverge faster than the  zeroth order piece as  $x \to 0^+$, and it would not be possible to treat the solution as an acceptable perturbation \cite{dg}.
Also, the choice $b_0=0$
selects perturbations with finite energy \footnote{Using the energy notion obtained by going to
second order perturbation theory \cite{ghi}. The one associated to ${\cal H}_z$ as a quantum Hamiltonian
is not relevant to the gravitational problem, in view of our comments above.}.

 Motivated by these facts, we choose as boundary condition at $x=0=r$
\begin{equation} \label{cero2}
 \Psi_z \simeq a_0 \left[ \left(\frac{x}{|M|}\right)^{1/2}  + ... \right] = a_0 \left[ \left(\frac{r}{|M|}\right)  + ... \right]
 \end{equation}
Our problem is now well defined: we are interested in solutions of (\ref{zerilli}) satisfying (\ref{crc}) and
(\ref{cero2}). Since ${\cal H}_z$ preserves these two conditions \cite{dg}, they will be preserved by the evolution
(\ref{zerilli}) if they are satisfied by the initial data $\Psi_z|_{t=0}$ and $\p \Psi_z / \p t |_{t=0}$.

An unstable solution of (\ref{zerilli}) for every harmonic number, satisfying (\ref{crc}) and (\ref{cero2}),  was found in \cite{dg}, 
and soon recognized by Cardoso and Cavaglia \cite{cardoso}  to correspond to Chandrasekhar's algebraic special modes \cite{prs}.
It is separable with  $E=-k^2$,
$\Psi_z^{unst} = \psi_{unst} e^{kt}, k = \frac{2\lambda (\lambda+1)}{6|M|}$ and
\begin{equation}
\label{ines}
 \psi_{unst}  = C_1 r{e^{-kr}} \frac{\left( r-2M \right) ^{2 k |M|}}{2 \lambda r + 6M} .
\end{equation}
The
existence of these modes  constitute
a proof of instability for the negative mass Schwarzschild spacetime.

There remains the question, however, of how
would one evolve arbitrary initial data satisfying (\ref{crc}) and (\ref{cero2}), since neither the space of functions nor
the differential operator in (\ref{zerilli}) fit in the standard settings of functional analysis and differential equation theory.
A change of   radial  variable and integrating factor is easily seen not to solve the problem, it  just moves the singularity
from the function to the integration measure. The tool we need is that of intertwining potentials, which uses
a map $\Psi_z \to \hat \Psi := \p \Psi_z  / \p x - g(x) \Psi_z$ that, with a suitable chosen $g$, transforms
a physically sensible $\Psi_z$ (i.e., satisfying (\ref{crc}) and (\ref{cero2})) into a smooth
and square integrable  $\hat \Psi$, obeying a Zerilli like equation with a potential $\hat V$ free
of singularities.
This provides  a method to evolve initial data, and to show that arbitrary initial data with compact support
away from $r=0$ does excite the
unstable modes (\ref{ines}).






\section{Intertwining operators} \label{inter}

Consider a two dimensional wave  equation with a space dependent potential $V$
\begin{equation} \label{z1}
\left[ \frac{\p^2 }{\p t^2} - \frac{\p^2 }{\p x^2} + V(x) \right] \Psi = 0,
\end{equation}
and a linear operator $\I= \frac{\p}{\p x} -g(x)$ such that \cite{price}
\begin{equation} \label{io}
\I \; \left[ - \frac{\p^2 }{\p x^2} + V(x) \right]  =   \left[ - \frac{\p^2 }{\p x^2} + \hat V(x) \right] \; \I
\end{equation}
for some potential $\hat V(x)$. Since $\I$  commutes with $\p /\p t$, any solution $\Psi$ of (\ref{z1})  gives a -possibly trivial- solution $\hat \Psi := \I \Psi$ for the equation
\begin{equation} \label{z2}
\left[ \frac{\p^2 }{\p t^2} - \frac{\p^2 }{\p x^2} + \hat V(x)
\right] \hat \Psi = 0.
\end{equation}
Separation of variables $\Psi = \exp(i\w t) \psi(x)$
  ( $\hat \Psi = \exp(i\w t) \hat \psi(x)$) reduces (\ref{z1})
and (\ref{z2})  to   Schr\"odinger like
equations
\begin{eqnarray} \label{z3}
{\cal H} \psi  &=& \left[ - \frac{\p^2 }{\p x^2} + V(x) \right] \psi = \w ^2 \psi ,\\
\hat {\cal H} \hat \psi &=& \left[ - \frac{\p^2 }{\p x^2} + \hat
V(x) \right] \hat \psi = \w ^2 \hat \psi , \label{z4}
\end{eqnarray}
If we do not specify boundary conditions, there will be
two linearly independent solutions of (\ref{z3}) for any chosen
eigenvalue. Let us   denote any two such solutions  as $\psi_{\w}^{(j)}, j=1,2$. Note from (\ref{io}) that
 $ \I \psi _{\w}^{(j)}$ are (possibly trivial) solutions of  (\ref{z4}).\\
The conditions for the existence of an intertwining operator can be
obtained
 by applying (\ref{io}) to an
arbitrary function $\psi$, and then isolating  terms in $\psi$ and $\psi' := \p \psi / \p x$.
The coefficient of $\psi'$ gives
\begin{equation} \label{vh}
\hat V = V - 2 g' ,
\end{equation}
and that of $\psi$ gives $(g' + g^2-V)' = 0$, i.e.
$g'+g^2=V-\w_o{}^2$ for some constant $\w_o$. This last condition is
more transparent in terms of  $\po =\exp (\int^x g(x') dx' )$, which satisfies $\po'/\po=g$:
\begin{equation} \label{zz}
\left[ - \frac{\p^2 }{\p x^2} + V \right] \po = \w_o ^2 \po,
\end{equation}
i.e.,  $\po$ is an $\w_o^2$ eigenvalue of ${\cal H}$. From this follows \cite{price},

\vspace{.3cm}

\noindent
{\bf Lemma 1:} From any eigenfunction (\ref{zz}) it is possible to construct an intertwining operator
 $\I= \frac{\p}{\p x} -g(x)$  by choosing
$g = \po' / \po$. This gives  $\hat V = V - 2 g'$ in (\ref{io}).

\vspace{.3cm}

Lemma 1 collects the results we  need from \cite{price},
but  we need to elaborate further on these results to get some information  about
the possible ways to invert the effect of $\I$.  To fix the notation, let  $\psi_{\w_o}^{(j=1)}=\po$,
$\psi_{\w_o}^{(j=2)}$ be linearly independent $\w_o^2$ eigenfunction of ${\cal H}$. The
kernel of $\I$ is the span of $\pe$, since   $0 = \I \psi = \psi' - \pe '
/ \pe \psi$ implies that $\psi$ is proportional to $\pe$. The form
of an intertwiner $\hat \I = \frac{\p}{\p x} -h(x)$
satisfying
\begin{equation} \label{iio}
\hat \I \; \left[ - \frac{\p^2 }{\p x^2} + \hat  V(x) \right]  =
\left[ - \frac{\p^2 }{\p x^2} +  V(x) \right] \; \hat \I
\end{equation}
can be guessed from Lemma 1 by noting that, since $\hat V - 2h' = V =
\hat V + 2g'$, the only possible way back to $V$ is to find an $\hat
{\cal H}$ eigenfunction $\hat \psi_o$ for which $\hat \psi_o ' / \hat \psi_o =-g$.
Since $g = \pe'/\pe$ this could only be possible if $ 1/\pe
$ were an eigenfunction of $\hat  {\cal H}$. That this is actually the
case can be checked by a direct calculation
 using our previous results, from where we obtain
  $\hat {\cal H}  (1/ \pe) = \w_o{}^ 2 /\pe$. We will set $\hat \psi {}_{\w_o}^{(j=2)} := 1/\pe$ and
  choose $\hat \psi {}_{\w_o}^{(j=1)}$ such that $\hat I \hat \psi {}_{\w_o}^{(j=1)} = \pe$.
  It follows that $\hat \I := \frac{\p}{\p x} + g(x)$ satisfies (\ref{iio}), and a simple calculation shows that
$\hat \I \I \psi = (\w_o{}^2 - {\cal H}) \psi$, i.e., the non trivial kernels of $\I$ and $\hat \I$ combine
in such a way that the kernel of $\hat \I \I$ is the two dimensional $\w_o^2$ eigenspace of ${\cal H}$. Note that we have
shown that we can label the $\w_o^2$ eigenfunctions
 such that
\begin{eqnarray} \nonumber
\pe &=&  \hat I \hat \psi _{\w_o}^{(j=1)}\\ \label{qsy}
I \psi _{\w_o}^{(j=2)} &=&  \hat \psi _{\w_o}^{(j=2)} = 1 / \pe \\
I \psi _{\w_o}^{(j=1)} &=& \hat I \hat \psi _{\w_o}^{(j=2)} =0\nonumber
\end{eqnarray}

\vspace{.3cm}

\noindent
{\bf Lemma 2:}  The kernel of $\I = \frac{\p}{\p x} - \psi _{\w_o}^{(j=1)} {}' / \psi _{\w_o}^{(j=1)}$
is the subspace spanned by $\po ^{(j=1)}$.  If $ \hat \I  =  \p/\p x +
 \psi _{\w_o}^{(j=1)} {}' / \psi _{\w_o}^{(j=1)}$, then (\ref{iio}) holds,
 $ \hat \I \I = (\w_o^2 - {\cal H})$, and  the ${\cal H}$ and $\hat {\cal H}$ eigenfunctions
can be labeled such that (\ref{qsy}) hold.

\vspace{.3cm}

The intertwining  operator (\ref{io}) will be useful whenever  $\hat V$ is simpler than
$V$. However, information is lost when solving (\ref{z2}) instead of (\ref{z1}), and we need
to know how to recover it. This problem is addressed in the theorem below.

\vspace{.3cm}

\noindent
{\bf Theorem:} Assume $\Psi(t,x)$ satisfies the wave equation (\ref{z1}) with initial conditions $\Psi(0,x)=f(x)$ and $\p \Psi / \p t (0,x)=
k(x)$. Let $\hat \Psi := \I \Psi, \hat f := \I f$ and $\hat k := \I k$, then:\\
(i) $\hat \Psi$ satisfies the wave equation (\ref{z2}) with
initial conditions $\hat \Psi(0,x)=\hat f(x)$ and $\p \hat \Psi / \p t (0,x)= \hat k(x)$. \\
(ii) If $\tilde \Psi := \hat \I \hat \Psi = \hat \I \I \Psi$
then
\begin{equation} \label{sol}
\Psi(t,x) =  \cos (\w_o t) f(x) + \frac{\sin (\w_o t)}{\w_o} k(x) + \frac{1}{w_o} \left(
\sin (\w_o t) \int_0^t \cos (\w_o t') \tilde \Psi(t',x) dt' -\cos (\w_o t) \int_0^t \sin (\w_o t') \tilde \Psi(t',x) dt' \right)
\end{equation}
unless  $\w_o = 0$, in which case
\begin{equation}  \label{sol2}
\Psi(t,x) = \int_0^t \left( \int_0^{t'} \tilde \Psi(t'',x) dt'' \right) dt' + t k(x) + f(x)
\end{equation}

\vspace{.3cm}

{\bf Proof:} (i) is trivial. To prove (ii) note from Lemma 2 that
\begin{equation} \label{eqn}
\tilde \Psi =  \hat \I \I \Psi =  (\w_o^2 - {\cal H}) \Psi = (\w_o^2 + \p^2  / \p t^2) \Psi
\end{equation}
where we have used that $\Psi$ satisfies  (\ref{z1}) in the last equality.
 The solution of (\ref{eqn}), regarded as a  differential equation on $\Psi$, is
\begin{equation} \label{sola}
\Psi(t,x) =  \cos (\w_o t) F(x) + \sin (\w_o t) K(x) + \frac{1}{w_o} \left(
\sin (\w_o t) \int_0^t \cos (\w_o t') \tilde \Psi(t',x) dt' -\cos (\w_o t) \int_0^t \sin (\w_o t') \tilde \Psi(t',x) dt' \right)
\end{equation}
if $\w_0^2 \neq 0$, and
\begin{equation}  \label{sol2a}
\Psi(t,x) = \int_0^t \left( \int_0^{t'} \tilde \Psi(t'',x) dt'' \right) dt' + t R(x) + Q(x)
\end{equation}
if $\w_o =0$. The unknown functions $F$ and $K$ ($Q$ and $R$) is the information about $\Psi$ that we have lost
when applying $\hat \I \I$. Fortunately, this information  is contained in the
initial conditions, since it can readily be seen that $F(x) = \Psi(0,x)= f(x) $ and $\w_o K(x)=\p \Psi / \p t (0,x) = k(x)$.
This gives (\ref{sol}), and (\ref{sol2}) follows similarly $\Box$\\

The theorem above provides a procedure  to solve eq.(\ref{z1}) given  $\Psi|_{(t=0)}=f(x)$ and $\p \Psi / \p t |_{(t=0)}=
k(x)$, using intertwining operators. The steps to follow  are:
\begin{enumerate}
\item  Find $\hat f := \I f$ and $\hat k := \I k$.
\item  Solve eq.(\ref{z2}) subject to $\hat \Psi |_{(t=0)} = \hat f(x)$ and $ \p \hat \Psi/ \p t |_{(t=0)} = \hat k (x)$.
This can be done by
using a mode expansion of $\hat {\cal H}$,
that is, for every eigenvalue $E$ of $\hat {\cal H}$ with normalized eigenfunction $\hat \psi_E$ ($\hat {\cal H} \hat \psi_E
 = E \hat \psi_E$),
solve the elementary ordinary differential equation
$\ddot a_E = -E a_E $ subject to the initial conditions $\dot a_E(0) = \int \hat \psi_E ^*   \hat k
   \; dx $ and $a_E(0) = \int \hat \psi_E ^* \hat f  \; dx$, then
   \begin{equation} \label{step}
 \hat \Psi(t,x) = \int dE a_E(t) \hat \psi_E(x) ,
 \end{equation}
where the integral above is symbolic, and is meant to denote an integral over the continuum piece
of the spectrum plus a sum over the discrete part of it.
\item Define $\tilde \Psi = \hat \I \hat \Psi$, where $\hat \Psi$ is given in (\ref{step}).
\item $\Psi(t,x)$ is given by eqn.(\ref{sol}) or (\ref{sol2}).
\end{enumerate}
Note that the initial conditions $f$ and $k$ have to be used twice because of the information
lost caused by applying $\I$ and $\hat I$, both of which have a nonzero kernel.


\section{Intertwining operator for the negative mass Zerilli equation} \label{apli}

The intertwining operator technique above can be used to reformulate the problem (\ref{zerilli})-(\ref{cero})-(\ref{crc})
and end up with a self adjoint operator in a Hilbert space. This could have been  guessed from Lemma 1: since generic
solutions of ${\cal H}_z \psi_E = E \psi_E$ behave as (\ref{crc}) \cite{dg}, there is a good chance
that the transformed potential (\ref{vh}) be nonsingular at $r_s$, and this may well be a consequence of
\begin{equation}\label{hpsi1}
\hat \Psi = \Psi_z' - g \Psi_z
\end{equation}
being a smooth function of the perturbed metric. All these expectations
turn out to be right. The ${\cal H}_z$ eigenfunction that we will use in Lemma 1 is the generalization to
arbitrary harmonic number $\ell$ of the zero mode ($\w_0=0$)
in \cite{ghi}. The zero mode equation is of the form,
\begin{equation}\label{zm1}
{\cal H}_z \psi_0 = 0
\end{equation}
with $V$ of the form (\ref{poten1}), and $M < 0$, and the resulting potential is then
\begin{equation} \label{zm4}
\hat V = V - 2 (\psi_0' / \psi_0 )'
\end{equation}
We need to find solutions $\psi_0$ such that $\hat V$ has no ``kinematic'' singularities. These may arise from singularities in $V$ or in $(\psi_0' / \psi_0 )'$. Let us first consider the singularity for $r=r_s$, with $r_s =-(3 M/\lambda)$. Using the fact that $\psi_0$ is a solution of (\ref{zm1}), and turning to $r$ (instead of $x$) derivatives, we find,
\begin{equation}\label{zm5}
\hat{V} = \frac{2(r-2M)^2}{r^2 \psi_0^2}\left(\frac{d \psi_0}{dr}\right)^2 -V(r)
\end{equation}
Now, if $\psi(r)$ is {\em any} solution of (\ref{zm1}), then, near $r=r_s$, we have,
\begin{equation}\label{zm6}
\psi(r) = a_0 (r-r_s)^{-1}+\frac{a_0(\ell+2)^2(\ell-1)^2}{12 M(\ell^2+\ell+1)}
+a_3(r-r_s)^2+{\cal{O}}\left((r-r_s)^3\right)
\end{equation}
where $a_0$, and $a_3$ are arbitrary constants. Replacing in (\ref{zm5}), assuming $a_0 \neq 0$,
and expanding in powers of $(r-r_s)$, we find,
\begin{equation}\label{zm7}
\hat V = \frac{(\ell^2+\ell+2)(\ell+2)^3(\ell-1)^3}{216 M^2} +\left[\frac{(\ell^2+\ell+1)(\ell+2)^4(\ell-1)^4}{648 M^3}-\frac{4 (\ell^2+\ell+1)^2 a_3}{3 a_0}\right](r-r_s)+{\cal{O}}\left((r-r_s)^2\right)
\end{equation}
which shows that $\hat V$ is smooth for $r=r_s$, provided $a_0 \neq 0$. If $a_0=0$, $\hat V$ has a second order
pole at $r_s$. We consider therefore, $\psi_0$ of the form,
\begin{equation}\label{zm2}
\psi_0=\frac{\chi(r)}{6 M+r(\ell+2)(\ell-1)}
\end{equation}
with $\chi$ smooth in $ r \geq 0$. Replacing (\ref{zm2}) in (\ref{zm1}), we find that $\chi$ satisfies,
\begin{equation}\label{zm8}
\frac{d^2 \chi}{dr^2} +\frac{\left[6M^2+2 r \lambda(3M-r)\right]}{r(r-2M)(3M+\lambda r)}\frac{d \chi}{dr}
-\frac{\left[6M^2+2 r \lambda(3M+\lambda r)\right]}{r^2(r-2M)(3M+\lambda r)}\  \chi =0
\end{equation}
Using this result we may write (\ref{zm4}) as,
\begin{equation}\label{zm5a}
\hat{V} = \frac{2(r-2M)^2}{r^2 \chi^2}\left(\frac{d\chi}{dr}\right)^2
-\frac{4(r-2M)^2 \lambda}{r^2 (3M+\lambda r) \chi}\frac{d\chi}{dr}-
\frac{(r-2M)(6M^2+\lambda r (2\lambda r+4M))} {r^4 (3M+\lambda r)}
\end{equation}

Then, since, as shown, $\hat V$ is actually free from the singularity for $r=r_s$, the only remaining possible singularities for $r >0$ would correspond to the zeros of $\chi$. Therefore, we must require that $\chi$ in (\ref{zm2}) be not only smooth, but also non vanishing for $r>0$. It turns out that (\ref{zm8}) admits, for every $\ell \geq 2$, a polynomial solution of the form,
\begin{equation} \label{zm3}
\chi(r)=\sum_{n=1}^{\ell+2}\frac{(n-2)\left[(n-4)\ell(\ell+1)+n-1\right]\Gamma(\ell+n-1)
}{ 2^{n}\Gamma(n)^2\Gamma(\ell-n+3) (-M)^{n}}r^n
\end{equation}

For $\ell=2$ we have the mode found in \cite{ghi},
\begin{equation} \label{zm3a}
\chi(r)=-\frac{3r}{2M}+\frac{3r^3}{4M^3}+\frac{r^4}{4M^4}
\end{equation}
which is positive for $M<0$ and $r>0$. Similarly, for $M<0$, and $\ell \geq 3$ we have,
\begin{eqnarray} \label{zm3b}
\chi(r) & = & \frac{r}{|M|} \left[\frac{3}{2}-\frac{\ell(\ell+2)(\ell^2-1)}{32 } \frac{r^2}{|M|^2}+\frac{\ell(\ell+2)(\ell^2-1)}{96}\frac{ r^3}{|M|^3} \right.
\nonumber \\ & &\left. +
\frac{\ell(\ell^2+\ell+4)(\ell+3)(\ell^2-4)(\ell^2-1)}{6144}\frac{ r^4}{|M|^4}+\dots\right]
\end{eqnarray}
where all the remaining terms, indicated by dots, are non negative for $r>0$. Consider now the fourth degree polynomial given explicitly between the brackets in (\ref{zm3b}). Clearly, it is positive (equal to $3/2$) for $r=0$, and it is also positive for sufficiently large $r$. Therefore, it can only have a zero if its derivative vanishes at least at one point for $r>0$. One can check that for $r > 0$ there is only one root given by,
\begin{equation}\label{zm9}
r_0 = \frac{4|M|\left(\sqrt{6\ell(\ell^2-1)(\ell+2)-108}-6\right)}{(\ell-2)(\ell+3)(\ell^2+\ell+4)}
\end{equation}
which must, therefore, correspond to a minimum of the polynomial in $r > 0$.
Replacing $r=r_0$ in (\ref{zm3b}) we find,
\begin{equation}\label{zm10}
\chi(r_0) \geq \frac{16 (\rho-6)\left((\ell^3-\ell)(\ell+2)-18\right)\left((\ell^3-\ell)(\ell+2)(\rho-18)+288\right)}
{(\ell^2+\ell+4)^4(\ell-2)^4(\ell+3)^4}
\end{equation}
where $\rho=\sqrt{6\ell(\ell^2-1)(\ell+2)-108}$. The right hand side of (\ref{zm10}) is positive for $\ell \geq 3$, and, therefore, $\chi(r) > 0$  for $r >0$, and the potential $\hat V$ is smooth for $r > 0$, and free from ''kinematical singularities'', as intended. The explicit form of $\hat V$ as a function of $r$ is very complicated but, fortunately, it is not required for the rest of our analysis.

We may, nevertheless, obtain several features of $\hat V$, directly from  (\ref{zm5a}). First, since $\chi$ is a polynomial of degree $L+2$, from (\ref{zm5a}) we find that for large $r$,
\begin{equation}\label{zm11a}
\hat V = \frac{(L+2)(L+1)}{r^2} + {\cal{O}}(r^{-3})
\end{equation}
and, therefore, $\hat V > 0$, and $\hat V \to 0$, for sufficiently large $r$. Also,
from (\ref{zm3}), for $r \to 0$ we have, in general
\begin{equation}\label{zm11}
\hat V = 12 M^2 r^{-4} -2 M (\ell^2+\ell+3) r^{-3} +{\cal{O}}(r^{-2})
\end{equation}
From (\ref{tortu}),  for $M<0$ we have
\begin{equation}\label{zm12}
r = 2 |M|^{1/2} x^{1/2}+\frac{2}{3} x +\frac{1}{18 |M|} x^{3/2}+{\cal{O}}(x^{2}),
\end{equation}
which, inserted in (\ref{zm12})  gives
\begin{equation}\label{zm13}
\hat V = \frac{3}{4 x^2} +\frac{\ell^2+\ell-1}{4|M|^{1/2}x^{3/2}} -\frac{\ell(\ell+1)}{4 |M|x} +{\cal{O}}(x^{-1/2})
\end{equation}
and
thus the general local solution of the differential equation $\hat {\cal H}_z \hat \psi_E = E \hat \psi_E$, for $x \to 0$ is of the form,
\begin{equation}
\hat \psi_E = a_0 \left( x^{3/2} + ... \right)  + b_0 \left(  x^{-1/2} + ... \right)
\end{equation}
which is not square integrable near $x=0$ unless $b_0=0$, as anticipated in Section \ref{rw} (see the discussion below eq. (\ref{vcero})).
As can be checked using (\ref{hpsi1}), this last condition
corresponds precisely to (\ref{cero2}).  Moreover, (\ref{hpsi1}) and (\ref{crc}) imply that $\hat \psi = \I \psi_z$ admits a Taylor
expansion around $r=r_s$, which is, of course consistent with fact that $\hat V$ is smooth there. Of particular relevance is the transformed of (\ref{ines}). This is of the form,
\begin{equation}\label{hines}
\hat \psi _{unst}  = C_1 r{e^{-kr}} \frac{\left( r-2M \right) ^{2 k |M|}}{2 \lambda r + 6M} \left[\frac{\ell(\ell+2)(\ell^2-1)}{12M}+\frac{r-2M}{r^2} \left(1 -\frac{r}{ \chi}\frac{d \chi}{dr}\right)\right]
\end{equation}
which, in accordance with our previous discussion, is smooth for $r=r_s$, and vanishes as $r^3$ for $ r \to 0$. Since $\chi$ is polynomial that does not vanish in $r > 0$, the function between square brackets on the right in (\ref{hines}) is also smooth and bounded in the full interval $0 \leq r \leq \infty$, and the function $\hat \psi$ is a square integrable eigenfunction of $\hat {\cal H}_z$, which, therefore, has, at least, one bound state. This, by the way, implies that $\hat V$ must have a region where it takes negative values, as can be explicitly checked for particular values of $L$. Figure 1 shows a graph of $\hat V$ and $\hat \psi _{unst}$ for $\ell=2$ and $M=-2$.\\

\begin{figure}[h]
\centerline{\includegraphics[height=9cm,angle=-90]{f1}}
\caption{Modified potential $\hat V$ for $\ell =2$ and $M=-2$ (solid line), and unstable mode (dotted line) as functions of $x$.
This mode satisfies $\hat {\cal H} \hat \psi_{unst} = -\hat \psi_{unst}$, see(\ref{ines}) and (\ref{hines}) .}
\end{figure}

The intertwining transformation changed the original problem into that of finding the solutions
of (\ref{z2}) with the potential (\ref{zm4}), which behaves as $ \frac{3}{4x^2} + ...  $ as $x \to 0^+$. Since
$\hat {\cal H}_z$ is a self adjoint operator on this space, arbitrary functions in this space can be
expanded using a basis of eigenfunctions of $\hat {\cal H}_z$, among which is the unstable mode (\ref{hines}).
Thus, generic initial data will excite this mode and  will therefore grow exponentially in time,
and this will also be true for the corresponding solution of the Zerilli equation, as follows from (\ref{sol2}).


\section{Summary} \label{sum}


Gravitational perturbations propagating on a negative mass Schwarzschild background is a subtle
problem for two reasons. First, this space is not globally hyperbolic. As a consequence,
the perturbation equations can be reduced to a single $1+1$ wave equation on a
semi-infinite domain $x>0$, $(t,x)$ standard inertial coordinates of two dimensional Minkowski space,
$x=0$ the position of the singularity. This implies that a physically motivated choice of boundary conditions
at $x=0$ is required. There is a unique choice dictated simultaneously by two conditions \cite{ghi,dg}: (i) that
the linearized regime be valid in the whole domain, and, in particular, that the invariants made out of the Riemman tensor
behave such that their first order piece does not diverge faster than their zeroth order piece as the singularity is approached;
 (ii) that the energy of
the perturbation, as measured using the second order correction to the Einstein tensor \cite{ghi} be finite.\\
The second issue is that  Zerilli's function $\Psi_z$ is a singular function of the first order metric coefficients. As a consequence, the wave
equation it obeys has  a potential  with a ``kinematic'' singularity, and it is not clear how to evolve
initial data, since the usual approach of separation of variables leading to a well behaved quantum Hamiltonian operator
for the $x$ coordinate breaks down.\\
In this paper we have shown that this second problem can be overcome using the technique of intertwining potentials
given in \cite{price} and references therein. It is shown that the function $\hat \Psi = \Psi_z - \psi_0'/\psi_0 \Psi_z =: {\cal I}
\Psi_z$
obeys a wave equation with a smooth potential if $\psi_0$ is a particular zero mode of Zerilli's Hamiltonian. Moreover,
the new Hamiltonian (spatial piece of the modified wave equation), naturally selects the boundary
condition that is physically relevant.\\
We have also shown that, in spite of the fact that ${\cal I}$ has a non trivial kernel, it is possible to
evolve the perturbation equations using, at two different steps, the initial condition for the Zerilli function.
A straightforward application of this formalism allows us to show that the unstable mode
found in \cite{dg} can actually be excited by initial data compactly supported away from the singularity. This
closes a gap in our proof in \cite{dg} of instability of the negative mass Schwarzschild spacetime.


\section*{Acknowledgments}


This work was supported in part by grants from CONICET (Argentina)
and Universidad Nacional de C\'ordoba.  RJG and GD are supported by
CONICET.




\begin{thebibliography}{99}
\bibitem{ghi} Gibbons G W, Hartnoll D and  Ishibashi A 2005 {\em Prog.Theor.Phys.}
{\bf 113} 963-978, hep-th/0409307.
\bibitem{dg}
R.~J.~Gleiser and G.~Dotti,
  %``Instability of the negative mass Schwarzschild naked singularity,''
  Class.Quant.Grav.{\bf 23}, 5063 (2006) [arXiv:gr-qc/0604021]
\bibitem{zerilli} F. J.
Zerilli, Phys. Rev. {\bf D2}, 2141 (1970).
\bibitem{ki}
  H.~Kodama and A.~Ishibashi,
  %``Stability of higher-dimensional Schwarzschild black holes,''
  Prog.\ Theor.\ Phys.\  {\bf 110}, 901 (2003)
  [arXiv:hep-th/0305185];
  %``A master equation for gravitational perturbations of maximally  symmetric
  %black holes in higher dimensions,''
  Prog.\ Theor.\ Phys.\  {\bf 110}, 701 (2003)
  [arXiv:hep-th/0305147];
  %``Brane world cosmology: Gauge-invariant formalism for perturbation,''
  Phys.\ Rev.\  D {\bf 62}, 064022 (2000)
  [arXiv:hep-th/0004160];
  \bibitem{price}
  A.~Anderson and R.~H.~Price,
  %``Intertwining of the equations of black hole perturbations,''
  Phys.\ Rev.\  D {\bf 43}, 3147 (1991).
  \bibitem{rw}
T. Regge, J. Wheeler, Phys. Rev. {\bf 108}, 1063 (1957).
\bibitem{meetz}
K. Meetz, Il Nuovo Cimento {\bf 34} 690 (1964).
  \bibitem{reed} M. Reed and B. Simon, {\em  Methods of modern mathematical physics} (v. 2: Fourier
  analysis, self-adjointness), section X,
  Academic Press (1975).
  \bibitem{cardoso}
  V.~Cardoso and M.~Cavaglia,
  %``Stability of naked singularities and algebraically special modes,''
  arXiv:gr-qc/0604101 (to appear in Phys. Rev. D)
  \bibitem{prs}
  S.~Chandrasekhar
  Proc.\ Roy.\ Soc.\ Lond.\ A {\bf 392}, 1 (1983).
  



\end{thebibliography}

\end{document}


\bibitem{doglepu}
 G.Dotti, R. J. Gleiser, J. Pullin
 %Instability of charged and rotating naked singularities
Phys.Lett. {\bf B644} (2007) 289-293 [arXiv:gr-qc/0607052]
\bibitem{cuatro}
  G.~Dotti, R.~J.~Gleiser, I.~F.~Ranea-Sandoval and H.~Vucetich,
  %``Gravitational instabilities in Kerr space times,''
  arXiv:0805.4306 [gr-qc].
\bibitem{cardoso}
  V.~Cardoso and M.~Cavaglia,
  %``Stability of naked singularities and algebraically special modes,''
   Phys.Rev. {\bf D74},   024027 (2006)
  [arXiv:gr-qc/0604101]
\bibitem{prs}
  S.~Chandrasekhar
  Proc.\ Roy.\ Soc.\ Lond.\ A {\bf 392}, 1 (1983).

